% Part: first-order-logic
% Chapter: syntax-and-semantics
% Section: free-vars-sentences

\documentclass[../../../include/open-logic-section]{subfiles}

\begin{document}

\olfileid{fol}{syn}{fvs}

\olsection{\printtoken{P}{variable} Bebas lan \printtoken{P}{sentence}}

\begin{defn}[Kemunculan bebas sawijining !!a{variable}]
\ollabel{defn:free-occ}
Kemunculan \emph{bebas} sawijining !!a{variable} ing !!a{formula}
ditetepake kanthi induktif kaya mangkene:
\begin{enumerate}
\item \indcase*{!A}{$!A$ atomik}{kabeh kemunculan !!{variable} ing
  $\indfrm$ iku bebas.}

\tagitem{prvNot}{\indcase{!A}{\lnot !B}{kemunculan !!{variable}
  kang bebas ing $\indfrm$ persis padha karo kang ana ing $!B$.}}{}

\item \indcase{!A}{(!B \ast !C)}{kemunculan
  !!{variable} kang bebas ing $\indfrm$ yaiku kang ana ing $!B$
  bebarengan karo kang ana ing~$!C$.}

\tagitem{prvAll}{\indcase{!A}{\lforall[x][!B]}{kemunculan !!{variable}
  kang bebas ing $\indfrm$ yaiku kabeh kang ana ing~$!B$ kajaba
  kemunculan~$x$.}}{}

\tagitem{prvEx}{\indcase{!A}{\lexists[x][!B]}{kemunculan !!{variable}
  kang bebas ing $\indfrm$ yaiku kabeh kang ana ing~$!B$ kajaba
  kemunculan~$x$.}}{}
\end{enumerate}
\end{defn}

\begin{defn}[Variabel kaiket]
Sawijining kemunculan !!a{variable} ing formula~$!A$ iku
\emph{kaiket} yen ora bebas.
\end{defn}

\begin{prob}
Wenehana definisi induktif kanggo kemunculan variabel kang kaiket
manut pola \olref[fol][syn][fvs]{defn:free-occ}.
\end{prob}

\begin{defn}[Cakupan]
\iftag{prvAll}{Yen $\lforall[x][!B]$ iku kemunculan sawijining subformula
  ing formula~$!A$, mula kemunculan~$!B$ kang cocog ing~$!A$
  diarani \emph{cakupan} saka kemunculan~$\lforall[x]$ kang cocog.
  \iftag{prvEx}{Semono uga kanggo $\lexists[x]$.}{}}{Yen
  $\lexists[x][!B]$ iku kemunculan sawijining subformula ing
  formula~$!A$, mula kemunculan~$!B$ kang cocog ing~$!A$
  diarani \emph{cakupan} saka kemunculan~$\lexists[x]$ kang cocog.}

Yen $!B$ iku cakupan sawijining kemunculan kuantor
\iftag{prvAll}{$\lforall[x]$\iftag{prvEx}{ utawa
    $\lexists[x]$}{}}{$\lexists[x]$} ing~$!A$, mula kemunculan bebas
$x$ ing~$!B$ dadi kaiket ing \iftag{prvAll}{$\lforall[x][!B]$\iftag{prvEx}{ lan
$\lexists[x][!B]$}{}}{$\lexists[x][!B]$}. Kita kandha yen
kemunculan kasebut \emph{diiket dening}
kemunculan kuantor kang kasebut mau.
\end{defn}

\begin{ex}
Gatekna formula iki:
\[
\lexists[\Obj v_0][\underbrace{\Atom{\Obj A^2_0}{\Obj v_0,\Obj v_1}}_{!B}] 
\]
$!B$ nuduhake cakupan $\lexists[\Obj v_0]$.
Kuantor kasebut ngiket kemunculan $\Obj v_0$ ing $!B$, nanging
ora ngiket kemunculan $\Obj v_1$. Dadi, $\Obj v_1$ iku
variabel bebas ing kasus iki.


Saiki kita bisa ndeleng carane iki lumaku ing !!{formula}~$!A$
kang luwih rumit:
\[
\lforall[\Obj v_0][\underbrace{(\Atom{\Obj A^1_0}{\Obj v_0} \lif
    \Atom{\Obj A^2_0}{\Obj v_0, \Obj v_1})}_{!B}] \lif \lexists[\Obj
  v_1][\underbrace{(\Atom{\Obj A^2_1}{\Obj v_0, \Obj v_1} \lor \lforall[\Obj v_0][\overbrace{\lnot \Atom{\Obj A^1_1}{\Obj v_0}}^{!D}])}_{!C}]
\]
$!B$ iku cakupan saka $\lforall[\Obj v_0]$ kang kapisan, $!C$
iku cakupan saka $\lexists[\Obj v_1]$, lan $!D$ iku cakupan saka
$\lforall[\Obj v_0]$ kang kapindho. $\lforall[\Obj v_0]$ kang kapisan
ngiket kemunculan $\Obj v_0$ ing~$!B$, $\lexists[\Obj v_1]$ ngiket
kemunculan $\Obj v_1$ ing $!C$, lan $\lforall[\Obj v_0]$ kang kapindho
ngiket kemunculan $\Obj v_0$ ing~$!D$. Kemunculan $\Obj v_1$
kang kapisan lan kemunculan $\Obj v_0$ kang kaping papat iku bebas
ing~$!A$. Kemunculan $\Obj v_0$ kang pungkasan iku bebas ing $!D$,
nanging kaiket ing $!C$ lan~$!A$.
\end{ex}

\begin{defn}[Ukara]
!!^a{formula}~$!A$ iku \article{sentence} \emph{!!{sentence}} yen lan
mung yen ora ngemot kemunculan bebas saka !!{variable}s.
\end{defn}

% tambahna tuladha!

\end{document}
